\documentclass[10pt,a4paper]{amsart}
\usepackage[margin=19mm]{geometry}
\usepackage{amsmath,amssymb,mathtools}
\usepackage[scr=rsfs,cal=euler]{mathalfa}
\usepackage{xfrac}
\usepackage[unicode,hidelinks,bookmarks=false]{hyperref}
\newcommand{\ie}{i.e.\ }
\newcommand{\cf}{cf.\ }
\newcommand{\resp}{respectively\ }
\newcommand{\Subsection}{\subsection}
\providecommand{\nobreakdash}{\nobreak\hbox{-}}
\setlength{\emergencystretch}{2em}
\multlinegap0pt
\usepackage{amssymb}
\usepackage{mathtools}
\usepackage{bm}
\usepackage{xcolor}
\newtheorem{proposition}{Proposition}
\newcommand{\bC}{\mathbf{C}}
\newcommand{\bD}{\mathbf{D}}
\newcommand{\bF}{\mathbf{F}}
\newcommand{\bx}{\mathbf{x}}
\newcommand{\lmax}{\lambda_{\max}}
\title{A Moment-Based Eulerian Method for Variance-Based Finite-Time Lyapunov Exponent Computation in Stochastic Flows}
\author{Shingyu Leung}
\date{Source: arXiv:2606.19681v1 (CC BY 4.0)}
\begin{document}
\maketitle

\noindent\emph{Source and licence.} A Moment-Based Eulerian Method for Variance-Based Finite-Time Lyapunov Exponent Computation in Stochastic Flows. Version arXiv:2606.19681v1 (CC BY 4.0), distributed by arXiv under the Creative Commons Attribution 4.0 International licence, http://creativecommons.org/licenses/by/4.0/, as recorded in the arXiv licence metadata for this exact version and shown on the abstract page https://arxiv.org/abs/2606.19681v1. Full licence notice: \texttt{licenses/arxiv-2606.19681-CC-BY-4.0.txt}.\par\smallskip

\noindent\textit{Editorial scope.} Counted unit: the article's proposition prop:diffusion\_monotonicity (source0.tex lines 495--504). The article's complete original proof of it is delivered with it (source0.tex lines 506--515, one \texttt{proof} environment of 10 lines). No mathematics is written, repaired or re-derived: the statement and the proof are the article's own lines, in the article's own notation, and no retained line is substituted or re-flowed.  No further source-local context is needed: every cross-reference of every retained line resolves inside the two delivered spans.  Nothing was omitted from any retained span, so the package carries no omission record.  No retained line contains a citation, so the wrapper carries no bibliography and no bibliography entry.  No retained line contains a figure environment, an image-inclusion command or drawing code, so no image is delivered and the package declares no asset dependency.  Editorial formatting only, in addition: an AMS \texttt{amsart} preamble that restates the article's own declarations and macros used by the retained lines, namely the environments \texttt{proposition} on separate counters, together with the standard packages those lines need.  The retained environments print in this document as Proposition 1 in order of appearance, whereas the article prints the same items with the numbering of its own full document.  The complete native source of the article as distributed in the arXiv e-print for this exact version is bundled unchanged as \texttt{sources/203162/source0.tex}.
\begin{proposition}[Monotonicity with respect to diffusion]\label{prop:diffusion_monotonicity}
Let \(\bC_T^{(1)}(\bx_0)\) and \(\bC_T^{(2)}(\bx_0)\) be the covariance matrices generated by two diffusion tensors \(\bD_1\) and \(\bD_2\) along the same deterministic trajectory, with zero initial covariance.  If \(\bD_1(\bx(t),t)\preceq \bD_2(\bx(t),t)\) for \(0\leq t\leq T\), then
\[
    \bC_T^{(1)}(\bx_0)\preceq \bC_T^{(2)}(\bx_0),
    \qquad
    \lmax\!\left(\bC_T^{(1)}(\bx_0)\right)
    \leq
    \lmax\!\left(\bC_T^{(2)}(\bx_0)\right).
\]
\end{proposition}

\begin{proof}
Subtracting the two representation formulas gives
\[
    \bC_T^{(2)}-\bC_T^{(1)}
    =
    \int_0^T
    \bF(T,s)\left(\bD_2(\bx(s),s)-\bD_1(\bx(s),s)\right)\bF(T,s)^T\,ds .
\]
The integrand is positive semidefinite for every \(s\), hence the integral is positive semidefinite.  This gives the Loewner-order inequality.  The largest-eigenvalue inequality follows from monotonicity of \(\lmax\) on symmetric matrices under the Loewner order.
\end{proof}

\end{document}
